\documentclass[a4paper]{article}

\usepackage{cv}
\hypersetup{pdftitle={Yan Lin - Cover Letter}}

\begin{document}
\cvname{Yan Lin}
\cvtagline{Assistant Professor / Department of Computer Science / Aalborg University}

Email: \href{mailto:lyan@cs.aau.dk}{lyan@cs.aau.dk} \,/\,
Homepage: \href{https://www.yanlincs.com}{yanlincs.com} \,/\,
Phone: +45 9186-1833

{\color{rule}\rule{\linewidth}{0.4pt}}

\vspace{\entryskip}
\today

\vspace{\entryskip}
Dear Members of the Assessment Committee,

\vspace{\entryskip}
I am writing to apply for the [position] in the [department], [institution].
My name is Yan Lin, and I am currently an Assistant Professor in the Department of Computer Science at Aalborg University.
I completed my PhD in Computer Science at Beijing Jiaotong University in March 2024.
I first joined Aalborg University in September 2022 as a visiting PhD student, returned as a Postdoctoral Researcher in June 2024, and took up my current Assistant Professorship in June 2026.
I would welcome the opportunity to continue my academic career in this position.

\vspace{\entryskip}
My doctoral thesis, Self-supervised Learning of Spatiotemporal Trajectory Data, introduced self-supervised representation learning techniques for spatiotemporal trajectory data.
Since 2019, I have published 33 peer-reviewed papers in leading venues such as IEEE TKDE, SIGMOD, NeurIPS, KDD, WWW, AAAI, and IJCAI, including a NeurIPS 2025 oral presentation.
Of these, 17 are first, co-first, or lead student author papers, and my work has over 1{,}000 citations on Google Scholar with an h-index of 14.
My research centers on spatiotemporal data mining, developing AI models for tasks such as trajectory representation, trajectory recovery, and travel time estimation.

\vspace{\entryskip}
I also work on cross-disciplinary AI collaborations.
I am part of the Inverse Design of Materials Using Diffusion Probabilistic Models project under the Villum Synergy Programme, in collaboration with the Department of Chemistry and Bioscience at Aalborg University, which applies AI to the discovery and design of disordered materials.
I have served as principal investigator on research projects and contributed to funded proposals supported by the Villum Foundation and the National Natural Science Foundation of China.
In this position, I would pursue grants such as Villum Synergy, Villum Experiment, and DFF research projects to support my research and provide training opportunities for students.

\vspace{\entryskip}
I integrate my research into my teaching.
The AI Systems \& Infrastructure course I teach draws on the practical knowledge I gained through my research and project work, covering AI model deployment, scalable infrastructure, and production-ready AI systems.
I follow Aalborg University's Problem Based Learning model, prepare course materials from scratch as accessible blog-style literature, and supervise Bachelor's and Master's semester projects on applied AI.
Students have praised the course structure, content, and teaching quality, and described my supervision as helpful and accessible.
In this position, I would develop new courses and supervise thesis projects that connect research with practical applications.

\vspace{\entryskip}
I maintain active engagement in the academic community.
I serve as Secretary of the IEEE (Denmark Section) Computer Society, as a reviewer for journals such as IEEE TKDE, ACM TKDD, IEEE TNNLS, ACM TIST, IEEE TII, and IEEE TVT, and as a program committee member for major conferences including ICLR, NeurIPS, KDD, CVPR, ICCV, AAAI, IJCAI, WWW, and WACV.
In this position, I would also take on administrative responsibilities such as coordinating courses or serving on departmental committees.

\vspace{\entryskip}
Enclosed, please find my curriculum vitae, a complete list of publications, my research plan, and my teaching portfolio.
Letters of reference can be provided upon request by Prof. Christian S. Jensen and Prof. Bin Yang, with whom I have collaborated on research, Prof. Morten M. Smedskjaer from the Department of Chemistry and Bioscience, with whom I collaborate on the Villum Synergy Programme, and Dr. Andreas Aakerberg, with whom I work closely to teach and supervise student groups.
I would be happy to provide any additional information if needed.
Thank you for your consideration.

\vspace{\entryskip}
Sincerely,

\vspace{\entryskip}
Yan Lin

\end{document}
